\section{Impostazione del problema e linee di analisi}
Muovendosi dal quadro definitorio introdotto (vedi  Capitolo~\ref{ch:introduzione-contesto}) e dagli obiettivi dello stage (vedi Capitolo~\ref{ch:obiettivi}), questo capitolo assume come fonte primaria le tabelle sui profili di attraversamento delle regioni d’informazione volo, con particolare riferimento alla sottosezione \textsc{FLIGHT AIRSPACES} (vedi Sezione \ref{sec:flight_airspace}) e ai collegamenti semantici con \textsc{FLIGHTS} (vedi Sezione \ref{sec:flights}).

Per evitare ambiguità, fissiamo la notazione di base impiegata nelle metriche calcolate sulle sequenze di attraversamento per ciascuna regione.
\begin{description}
  \item[$f$] Identificativo del volo nel sistema \gls{adrr}.
  \item[$R$] Regione di volo attraversata.
  \item[$t^{\mathrm{in}}_{f,R}$] Istante di ingresso del volo \(f\) nella regione \(R\).
  \item[$t^{\mathrm{out}}_{f,R}$] Istante di uscita del volo \(f\) dalla regione \(R\).
\end{description}
In simboli, useremo le seguenti equivalenze notazionali:
\[
  f \equiv \texttt{ECTRL\_ID},\quad
  R \equiv \texttt{FIR ID},\quad
  t^{\mathrm{in}}_{f,R} \equiv \texttt{Entry Time},\quad
  t^{\mathrm{out}}_{f,R} \equiv \texttt{Exit Time}.
\]
\clearpage
\section{Modello temporale e granularità}
\label{sec:modello-temporale}
Per passare da una rappresentazione continua a una rappresentazione discreta del tempo, si discretizza l’orizzonte di analisi suddividendolo in intervalli elementari tutti della stessa durata. Tale durata, indicata con \(\lambda\), rappresenta la granularità temporale dell’analisi (cioè il passo minimo con cui vengono osservate le grandezze di interesse).
Fissiamo un orizzonte di analisi \([\tau_0,\tau_1)\subset\mathbb{R}\) e scegliamo un passo di discretizzazione \(\lambda>0\). Definiamo quindi la griglia \(\mathcal{T}\) come la successione di nodi equispaziati a partire da \(\tau_0\):
\begin{equation*}
\label{eq:griglia-def}
\mathcal{T}\coloneqq \bigl\{\, t_k = \tau_0 + k\,\lambda \;\bigm|\; k=0,1,\dots \,\bigr\}, \qquad t_k \in [\tau_0,\tau_1) .
\end{equation*}
I nodi \(t_k\) scandiscono l’intervallo di analisi con passo \(\lambda\).
Un esempio schematico della griglia così definita è mostrato in \figurename~\ref{fig:griglia-temporale}.
\begin{figure}[htbp]
\centering
\begin{tikzpicture}[x=1cm,y=1cm]
  % Parametri (solo per disegno)
  \def\tauZero{0}
  \def\tauOne{10}
  \def\DeltaLen{1.4}

  % Asse temporale
  \draw[timeline] (\tauZero-0.6,0) -- (\tauOne+0.6,0);

  % Estremi dell'orizzonte [tau0, tau1)
  \fill (\tauZero,0) circle (2pt) node[below=2pt] {$\tau_0$};
  \draw[open endpoint] (\tauOne,0) circle (2pt) node[below=2pt] {$\tau_1$};

  % Nodi t_k equispaziati: k = 0,1,...,K (qui K=6 per illustrazione)
  \foreach \k in {0,...,6}{
    \draw[tick] ({\tauZero + \k*\DeltaLen},-0.12) -- ({\tauZero + \k*\DeltaLen},0.12);
  }
  % Etichette: t_0, t_1, t_2, ..., t_K
  \node[nodelab] at (\tauZero, -0.45) {$t_0$};
  \node[nodelab] at ({\tauZero + 1*\DeltaLen}, -0.45) {$t_1$};
  \node[nodelab] at ({\tauZero + 2*\DeltaLen}, -0.45) {$t_2$};
  \node at ({\tauZero + 3.2*\DeltaLen}, -0.45) {$\cdots$};
  \node[nodelab] at ({\tauZero + 6*\DeltaLen}, -0.45) {$t_K$};

  % Braccetto che mostra \Delta tra due nodi consecutivi
  \draw[brace] ({\tauZero + 2*\DeltaLen},0.6) -- node[above=6pt] {$\lambda$} ({\tauZero + 3*\DeltaLen},0.6);

  % Braccio che evidenzia l'orizzonte [tau0, tau1)
  \draw[brace] (\tauZero,1.5) -- node[above=6pt] {$T=[\tau_0,\tau_1)$} (\tauOne,1.5);
\end{tikzpicture}
\caption{Griglia temporale su \([\tau_0,\tau_1)\). L’estremo destro è aperto.}
\label{fig:griglia-temporale}
\end{figure}

Osserviamo che la definizione di \(\mathcal{T}\) introduce, in prima istanza, una successione potenzialmente infinita di nodi \((t_k)_{k\in\mathbb{N}}\); nella pratica, però, all’interno dell’intervallo \([\tau_0,\tau_1)\) si utilizzano solo un numero finito di istanti. Per poter lavorare con vettori e somme indicizzati nel tempo è quindi utile rendere esplicita la cardinalità effettiva della griglia all’interno dell’orizzonte di analisi. A tal fine si introduce un indice terminale
\begin{equation}
\label{eq:define-K}
K \;=\; \Bigl\lfloor \frac{\tau_1 - \tau_0}{\lambda} \Bigr\rfloor .
\end{equation}
\subsection{Finestre temporali}\label{sec:finestre_temporali}
Dato un parametro di ampiezza \(W>0\), a ciascun nodo \(t_k\) associamo la finestra semi–aperta a destra:
\begin{equation*}
\label{eq:finestra-def}
I_k \coloneqq [\,t_k,\; t_k+W)\,.
\end{equation*}
Nel seguito assumiamo \(\lambda \le W\) così che le finestre risultino contigue o sovrapposte; il caso \(\lambda>W\) non è considerato. Le due configurazioni sono illustrate in \figurename~\ref{fig:finestre-varianti} per completezza.
\clearpage
\begin{figure}[H]
  \centering
  % --- (B) W = Δ ---
  \begin{subfigure}{\linewidth}
    \centering
    \begin{tikzpicture}[x=1cm,y=1cm]
      \def\tauZero{0}\def\tauOne{10}\def\DeltaLen{1.6}\def\h{0.6}\def\winfill{black!25}\def\Wlen{1.6}
      \draw[timeline] (\tauZero-0.6,0) -- (\tauOne+0.6,0);
      \foreach \k in {0,1,2,3,4,5}{
        \coordinate (tkB\k) at ({\tauZero + \k*\DeltaLen},0);
        \draw[tick] (tkB\k)++(0,-0.12) -- +(0,0.24);
        \node[nodelab, yshift=-6pt] at (tkB\k) {$t_{\k}$};
      }
      \foreach \k in {0,1,2,3,4}{
        \draw[draw=black, fill=\winfill, fill opacity=0.35, line width=0.6pt]
          (tkB\k)++(0,-\h/2) rectangle ++(\Wlen,\h);
        \node[closed endpoint] at (tkB\k) {};
        \node[open endpoint] at ($(tkB\k)+(\Wlen,0)$) {};
      }
      \draw[brace] ($(tkB2)+(0,0.4)$) -- node[above=15pt, nodelab] {$\lambda=W$} ($(tkB3)+(0,0.4)$);
    \end{tikzpicture}
    \caption{\(W=\lambda\)\,: finestre contigue.}
  \end{subfigure}

  \vspace{1.5em}

  % --- (C) W > Δ ---
  \begin{subfigure}{\linewidth}
    \centering
    \begin{tikzpicture}[x=1cm,y=1cm]
      \def\tauZero{0}\def\tauOne{10}\def\DeltaLen{1.6}\def\h{0.6}\def\winfill{black!25}\def\Wlen{2.2}\def\ysep{0.9}
      \draw[timeline] (\tauZero-0.6,0) -- (\tauOne+0.6,0);
      \foreach \k in {0,1,2,3,4,5}{
        \coordinate (tkC\k) at ({\tauZero + \k*\DeltaLen},0);
        \draw[tick] (tkC\k)++(0,-0.12) -- +(0,0.24);
        \node[nodelab] at (tkC\k) {$t_{\k}$};
      }
      \foreach \k in {0,1,2,3,4}{
        \pgfmathtruncatemacro{\sgn}{mod(\k,2)==0 ? 1 : -1}
        \coordinate (wkC\k) at ($(tkC\k)+(0,\sgn*\ysep)$);
        \draw[draw=black, fill=\winfill, fill opacity=0.35, line width=0.6pt]
          (wkC\k)++(0,-\h/2) rectangle ++(\Wlen,\h);
        \node[closed endpoint] at (wkC\k) {};
        \node[open endpoint] at ($(wkC\k)+(\Wlen,0)$) {};
      }
      \draw[brace] ($(tkC1)+(0,-0.45)$) -- node[below=5pt, nodelab] {$\Delta$} ($(tkC0)+(0,-0.45)$);
      \draw[brace] ($(wkC1)+(\Wlen,-\h/2-0.15)$) -- node[below=5pt, nodelab] {$W$} ($(wkC1)+(0,-\h/2-0.15)$);
    \end{tikzpicture}
    \caption{\(W>\lambda\)\,: finestre sovrapposte.}
  \end{subfigure}

  \caption{Configurazioni delle finestre in funzione di \(W\) e \(\lambda\).}
  \label{fig:finestre-varianti}
\end{figure}

\clearpage
\section{Metriche}
In questa sezione si introducono le grandezze utilizzate per descrivere il carico di traffico su una regione di spazio aereo \(R\).  
A partire dalla griglia temporale \(\mathcal{T}\) e dalla finestra \(I_k\), si definiscono due metriche: il \emph{throughput}, che misura il numero di ingressi in una regione nell’arco di una finestra, e l’\emph{occupancy}, che quantifica invece quanti velivoli risultano presenti nella regione per l’intera durata della finestra. Queste quantità consentono di passare dai dati puntuali di ingresso/uscita dei singoli voli a misure sintetiche, più adatte sia all’analisi statistica sia agli utilizzi dei capitoli successivi.
\label{sec:metriche}
\subsection{Throughput}
\label{subsec:throughput-teoria}
Sulla griglia temporale \(\mathcal{T}\) definita in Sezione \ref{sec:modello-temporale}, e con la finestre semi–aperta \(I_k=[t_k,t_k+W)\) introdotta in Sezione \ref{sec:finestre_temporali}, il \emph{throughput} per una regione \(R\) è il numero di istanti di ingresso che cadono in \(I_k\).
    Denotiamo tale conteggio con \(T_R(I_k)\):
    \begin{equation}
    \label{eq:throughput-def}
        T_R(I_k)\;\coloneqq\; \#\bigl\{\, f \in \mathcal{F}_R \;\bigm|\; t_k \le t^{\mathrm{in}}_{f,R} < t_k+W \,\bigr\}.
    \end{equation}
    In questo caso \(\mathcal{F}_R\) è l’insieme dei voli che attraversano \(R\) e \(t^{\mathrm{in}}_{f,R}\) è l’istante di ingresso del volo \(f\) in \(R\).
    La scelta della finestra semi–aperta \([t_k,\,t_k+W)\) evita doppi conteggi sugli estremi condivisi.

    
    \begin{figure}[H]
\centering

\tikzset{
  region/.style={draw, rounded corners=2pt, fill=gray!20},
  traj/.style={line width=0.5pt, -Latex},
  proj/.style={dashed},
  taxis/.style={-Latex, thick},
  tmark/in/.style={circle, fill=black, inner sep=1.2pt},
  tmark/out/.style={circle, draw=black, fill=white, inner sep=1.2pt},
  openmark/.style={circle, draw=black, fill=white, inner sep=1.1pt},
  rombo/.style={diamond, draw=black, fill=black, minimum size=6pt, inner sep=0pt},
  note/.style={draw, fill=white, rounded corners=2pt, font=\scriptsize, inner sep=2pt}
}

\def\yaxis{-1.0}
\def\tk{1.0}
\def\tkW{4.0}

\begin{subfigure}[t]{.48\textwidth}
\centering
\begin{tikzpicture}[x=0.9cm,y=0.9cm]
  \def\yaxis{-1.2}
  \def\tinA{2.2}
  \draw[region] (0,0) rectangle (6,3);
  \node[font=\small] at (6,3.5) {$R$};
  \draw[traj] (-0.6,2.4) .. controls (1.2,2.1) and (3.2,1.8) .. (6.6,1.5);
  %\node[rombo] at (2.2,1.95) {};
  \fill[gray!20, draw=black!50] (\tk,\yaxis+0.22) rectangle (\tkW,\yaxis-0.22);
  \coordinate (EinA) at (0,2.3);
  \fill[red] (EinA) circle (1.6pt);
  \node[font=\scriptsize, anchor=south west] at (EinA) {$t^{\mathrm{in}}_{R}=2.2$};
  \draw[proj] (EinA) -- (\tinA,\yaxis);
  \draw[taxis] (-0.2,\yaxis) -- (6.6,\yaxis);
  \foreach \x in {0,...,6}{
    \draw (\x,\yaxis+0.08) -- (\x,\yaxis-0.08);
    \node[font=\scriptsize] at (\x,\yaxis-0.35) {\x.0};
  }
  \node[tmark/in] at (\tinA,\yaxis) {};
\end{tikzpicture}
\subcaption{Rotta A: ingresso \textbf{dentro} la finestra.}
\end{subfigure}
\hfill
\begin{subfigure}[t]{.48\textwidth}
\centering
\begin{tikzpicture}[x=0.9cm,y=0.9cm]
  \def\yaxis{-1.2}
  \def\tinB{4.8}
  \draw[region] (0,0) rectangle (6,3);
  \node[font=\small] at (6,3.5) {$R$};
  \draw[traj] (-0.6,0.7) .. controls (1.0,1.3) and (3.5,2.5) .. (6.6,2.8);
  %\node[rombo] at (4.8,2.2) {};
  \fill[gray!20, draw=black!50] (\tk,\yaxis+0.22) rectangle (\tkW,\yaxis-0.22);
  \coordinate (EinB) at (0,0.9);
  \fill[red] (EinB) circle (1.6pt);
  \node[font=\scriptsize, anchor=south west, yshift=0.5cm] at (EinB) {$t^{\mathrm{in}}_{R}=4.8$};
  \draw[proj] (EinB) -- (\tinB,\yaxis);
  \draw[taxis] (-0.2,\yaxis) -- (6.6,\yaxis);
  \foreach \x in {0,...,6}{
    \draw (\x,\yaxis+0.08) -- (\x,\yaxis-0.08);
    \node[font=\scriptsize] at (\x,\yaxis-0.35) {\x.0};
  }
  \node[tmark/in] at (\tinB,\yaxis) {};
\end{tikzpicture}
\subcaption{Rotta B: ingresso \textbf{fuori} dalla finestra.}
\end{subfigure}

\caption{Throughput $T_R(I_k)$ per la regione $R$ come conteggio degli ingressi che cadono nella finestra $[t_k,t_k{+}W)$.}
\label{fig:throughput-subfig}
\end{figure}

    
    \noindent La definizione \eqref{eq:throughput-def} è equivalente a un calcolo basato sulla differenza di due cumulate degli ingressi. Introduciamo la cumulata stretta $N_R(t)$ che conta gli ingressi. La proposizione seguente formalizza l’identità che adottiamo nell’implementazione.
    
    \begin{proposition}[Identità a due cumulate]
    \label{prop:two-cumulatives}
    Sia $I_k=[t_k,t_k+W)$ con $W>0$ e sia
    \[
    N_R(t_k)\;\coloneqq\;\#\bigl\{\,f\in\mathcal{F}_R:\ t^{\mathrm{in}}_{f,R}<t_k\,\bigr\}
    \]
    la cumulata stretta degli ingressi in $R$. Si avrà dunque che
    \begin{equation}
    \label{eq:throughput-diff-cumulate}
    T_R(I_k)\;=\;N_R(t_k{+}W)\;-\;N_R(t_k)\quad\forall{k}>0.
    \end{equation}
    \end{proposition}
    
    \begin{proof}
    Definiamo $S_R(t)=\{\,f\in\mathcal{F}_R:\ t^{\mathrm{in}}_{f,R}<t\,\}$, cosicchè $N_R(t)=\#S_R(t)$.
    Per $I_k=[t_k,t_k{+}W)$ l’insieme degli ingressi che cadono nella finestra è
    \[
    \begin{aligned}
    A_k
    &= \bigl\{\,f:\ t_k \le t^{\mathrm{in}}_{f,R} < t_k+W\,\bigr\} 
    \\[2pt]
    &= \bigl\{\,f:\ (t^{\mathrm{in}}_{f,R} < t_k+W)\ \wedge\ \neg(t^{\mathrm{in}}_{f,R} < t_k)\,\bigr\}
    \\[2pt]
    &= \bigl\{\,f:\ t^{\mathrm{in}}_{f,R} < t_k+W\,\bigr\}
       \ \cap\ 
       \bigl\{\,f:\ \neg(t^{\mathrm{in}}_{f,R} < t_k)\,\bigr\}
    \\[2pt]
    &= \bigl\{\,f:\ t^{\mathrm{in}}_{f,R} < t_k+W\,\bigr\}
       \ \setminus\ 
       \bigl\{\,f:\ t^{\mathrm{in}}_{f,R} < t_k\,\bigr\}
    \\[2pt]
    &= S_R(t_k{+}W)\ \setminus\ S_R(t_k).
    \end{aligned}
    \]
    Poiché $W>0$ implica $S_R(t_k)\subseteq S_R(t_k{+}W)$, passando alle cardinalità otteniamo
    \[
    \begin{aligned}
    \#A_k
    &=\#S_R(t_k{+}W)-\#S_R(t_k)
    \\[2pt]
    &=N_R(t_k{+}W)-N_R(t_k) \\[2pt]
    &=T_R(I_k)
    \end{aligned}
    \]
    che è proprio~\eqref{eq:throughput-diff-cumulate}.
    \end{proof}
    
    
    
    
\subsection{Occupancy}
\label{subsec:occupancy-teoria}
Sulla griglia temporale \(\mathcal{T}\) definita in \S\ref{sec:modello-temporale}, e con le finestre semi–aperte \(I_k=[t_k,t_k+W)\) introdotte in \S\ref{sec:finestre_temporali}, definiamo l’\emph{occupancy} come il numero di velivoli che risultano presenti in \(R\) per tutta la finestra:
    \begin{equation}
    \label{eq:occ-finestra-copertura}
    O_R(k) \;\coloneqq\; 
    \#\Bigl\{\, f \in \mathcal{F}_R \;\Bigm|\; 
    t^{\mathrm{in}}_{f,R} \le t_k \ \ \wedge\ \ t^{\mathrm{out}}_{f,R} \ge t_k{+}W \Bigr\}.
    \end{equation}
    In questo caso \(\mathcal{F}_R\) è l’insieme dei voli che attraversano \(R\), \(t^{\mathrm{in}}_{f,R}\) è l’istante di ingresso del volo \(f\) in \(R\) e \(t^{\mathrm{out}}_{f,R}\) è l’istante di uscita del volo \(f\) in \(R\).
    \clearpage
    \begin{figure}[H]
\centering

\tikzset{
  region/.style={draw, rounded corners=2pt, fill=gray!20},
  traj/.style={line width=0.5pt, -Latex},
  proj/.style={dashed},
  taxis/.style={-Latex, thick},
  tmark/in/.style={circle, fill=black, inner sep=1.2pt},
  tmark/out/.style={circle, draw=black, fill=white, inner sep=1.2pt},
  openmark/.style={circle, draw=black, fill=white, inner sep=1.1pt},
  rombo/.style={diamond, draw=black, fill=black, minimum size=6pt, inner sep=0pt},
  note/.style={draw, fill=white, rounded corners=2pt, font=\scriptsize, inner sep=2pt}
}

\def\yaxis{-1.0}
\def\tk{1.0}
\def\tkW{4.0}

\begin{subfigure}[t]{.48\textwidth}
\centering
\begin{tikzpicture}[x=0.9cm,y=0.9cm]
  \def\yaxis{-1.2}
  \def\tinA{0.8}
  \def\toutA{4.8}
  \draw[region] (0,0) rectangle (6,3);
  \node[font=\small] at (6,3.5) {$R$};

  % Traiettoria (passa dentro R)
  \draw[traj] (-0.6,2.2) .. controls (1.0,2.0) and (3.5,1.6) .. (6.6,1.4);

  % Finestra sulla linea temporale
  \fill[gray!20, draw=black!50] (\tk,\yaxis+0.22) rectangle (\tkW,\yaxis-0.22);

  % Punti di ingresso/uscita su frontiera R
  \coordinate (EinA)  at (0,2.15);
  \coordinate (EoutA) at (6,1.45);
  \fill[red] (EinA) circle (1.6pt);                       % ingresso pieno
  \node[openmark] at (EoutA) {};                           % uscita "open"
  \node[font=\scriptsize, anchor=south west] at (EinA)  {$t^{\mathrm{in}}_{R}=0.8$};
  \node[font=\scriptsize, anchor=south east] at (EoutA) {$t^{\mathrm{out}}_{R}=4.8$};

  % Proiezioni su asse temporale
  \draw[proj] (EinA) -- (\tinA,\yaxis);
  \draw[proj] (EoutA) -- (\toutA,\yaxis);

  % Asse temporale
  \draw[taxis] (-0.2,\yaxis) -- (6.6,\yaxis);
  \foreach \x in {0,...,6}{
    \draw (\x,\yaxis+0.08) -- (\x,\yaxis-0.08);
    \node[font=\scriptsize] at (\x,\yaxis-0.35) {\x.0};
  }

  % Marker ingresso/uscita su asse
  \node[tmark/in]  at (\tinA,\yaxis) {};
  \node[tmark/out] at (\toutA,\yaxis) {};
\end{tikzpicture}
\subcaption{Rotta A: \textbf{copre interamente} la finestra}
\end{subfigure}
\hfill
\begin{subfigure}[t]{.48\textwidth}
\centering
\begin{tikzpicture}[x=0.9cm,y=0.9cm]
  \def\yaxis{-1.2}
  \def\tinB{1.4}
  \def\toutB{3.2}
  \draw[region] (0,0) rectangle (6,3);
  \node[font=\small] at (6,3.5) {$R$};

  % Traiettoria
  \draw[traj] (-0.6,0.9) .. controls (1.2,1.4) and (3.4,2.4) .. (6.6,2.7);

  % Finestra sulla linea temporale
  \fill[gray!20, draw=black!50] (\tk,\yaxis+0.22) rectangle (\tkW,\yaxis-0.22);

  % Punti di ingresso/uscita su frontiera R
  \coordinate (EinB)  at (0,1.05);
  \coordinate (EoutB) at (6,2.62);
  \fill[red] (EinB) circle (1.6pt);
  \node[openmark] at (EoutB) {};
  \node[font=\scriptsize, anchor=south west, yshift=1em] at (EinB)  {$t^{\mathrm{in}}_{R}=1.4$};
  \node[font=\scriptsize, anchor=south east, yshift=-0.4em] at (EoutB) {$t^{\mathrm{out}}_{R}=3.2$};

  % Proiezioni su asse temporale
  \draw[proj] (EinB)  -- (\tinB,\yaxis);
  \draw[proj] (EoutB) -- (\toutB,\yaxis);

  % Asse temporale
  \draw[taxis] (-0.2,\yaxis) -- (6.6,\yaxis);
  \foreach \x in {0,...,6}{
    \draw (\x,\yaxis+0.08) -- (\x,\yaxis-0.08);
    \node[font=\scriptsize] at (\x,\yaxis-0.35) {\x.0};
  }

  % Marker ingresso/uscita su asse
  \node[tmark/in]  at (\tinB,\yaxis) {};
  \node[tmark/out] at (\toutB,\yaxis) {};
  
\end{tikzpicture}
\subcaption{Rotta B: \textbf{non copre} la finestra}
\end{subfigure}

\caption{Occupazione \(O_R(k)\) per la regione \(R\): conteggio dei voli la cui permanenza copre interamente la finestra \(I_k=[t_k,t_k{+}W)\).}
\label{fig:occupazione-subfig}
\end{figure}
    
    L'implementazione algoritmica non valuta la condizione di copertura \eqref{eq:occ-finestra-copertura} per ogni volo e ogni finestra, ma sfrutta il principio dell'array di differenze. La proposizione seguente formalizza l’identità che adottiamo nell’implementazione.
    
    \subsection{Occupancy}
\label{subsec:occupancy-teoria}
Sulla griglia temporale \(\mathcal{T}\) definita in \S\ref{sec:modello-temporale}, e con le finestre semi–aperte \(I_k=[t_k,t_k+W)\) introdotte in \S\ref{sec:finestre_temporali}, definiamo l’\emph{occupancy} come il numero di velivoli che risultano presenti in \(R\) per tutta la finestra:
\begin{equation}
\label{eq:occ-finestra-copertura}
O_R(k) \;\coloneqq\; 
\#\Bigl\{\, f \in \mathcal{F}_R \;\Bigm|\; 
t^{\mathrm{in}}_{f,R} \le t_k \ \ \wedge\ \ t^{\mathrm{out}}_{f,R} \ge t_k{+}W \Bigr\}.
\end{equation}
In questo caso \(\mathcal{F}_R\) è l’insieme dei voli che attraversano \(R\), \(t^{\mathrm{in}}_{f,R}\) è l’istante di ingresso del volo \(f\) in \(R\) e \(t^{\mathrm{out}}_{f,R}\) è l’istante di uscita del volo \(f\) in \(R\).

\clearpage
\begin{figure}[H]
\centering

\tikzset{
  region/.style={draw, rounded corners=2pt, fill=gray!20},
  traj/.style={line width=0.5pt, -Latex},
  proj/.style={dashed},
  taxis/.style={-Latex, thick},
  tmark/in/.style={circle, fill=black, inner sep=1.2pt},
  tmark/out/.style={circle, draw=black, fill=white, inner sep=1.2pt},
  openmark/.style={circle, draw=black, fill=white, inner sep=1.1pt},
  rombo/.style={diamond, draw=black, fill=black, minimum size=6pt, inner sep=0pt},
  note/.style={draw, fill=white, rounded corners=2pt, font=\scriptsize, inner sep=2pt}
}

\def\yaxis{-1.0}
\def\tk{1.0}
\def\tkW{4.0}

\begin{subfigure}[t]{.48\textwidth}
\centering
\begin{tikzpicture}[x=0.9cm,y=0.9cm]
  \def\yaxis{-1.2}
  \def\tinA{0.8}
  \def\toutA{4.8}
  \draw[region] (0,0) rectangle (6,3);
  \node[font=\small] at (6,3.5) {$R$};

  % Traiettoria (passa dentro R)
  \draw[traj] (-0.6,2.2) .. controls (1.0,2.0) and (3.5,1.6) .. (6.6,1.4);

  % Finestra sulla linea temporale
  \fill[gray!20, draw=black!50] (\tk,\yaxis+0.22) rectangle (\tkW,\yaxis-0.22);

  % Punti di ingresso/uscita su frontiera R
  \coordinate (EinA)  at (0,2.15);
  \coordinate (EoutA) at (6,1.45);
  \fill[red] (EinA) circle (1.6pt);                       % ingresso pieno
  \node[openmark] at (EoutA) {};                           % uscita "open"
  \node[font=\scriptsize, anchor=south west] at (EinA)  {$t^{\mathrm{in}}_{R}=0.8$};
  \node[font=\scriptsize, anchor=south east] at (EoutA) {$t^{\mathrm{out}}_{R}=4.8$};

  % Proiezioni su asse temporale
  \draw[proj] (EinA) -- (\tinA,\yaxis);
  \draw[proj] (EoutA) -- (\toutA,\yaxis);

  % Asse temporale
  \draw[taxis] (-0.2,\yaxis) -- (6.6,\yaxis);
  \foreach \x in {0,...,6}{
    \draw (\x,\yaxis+0.08) -- (\x,\yaxis-0.08);
    \node[font=\scriptsize] at (\x,\yaxis-0.35) {\x.0};
  }

  % Marker ingresso/uscita su asse
  \node[tmark/in]  at (\tinA,\yaxis) {};
  \node[tmark/out] at (\toutA,\yaxis) {};
\end{tikzpicture}
\subcaption{Rotta A: \textbf{copre interamente} la finestra}
\end{subfigure}
\hfill
\begin{subfigure}[t]{.48\textwidth}
\centering
\begin{tikzpicture}[x=0.9cm,y=0.9cm]
  \def\yaxis{-1.2}
  \def\tinB{1.4}
  \def\toutB{3.2}
  \draw[region] (0,0) rectangle (6,3);
  \node[font=\small] at (6,3.5) {$R$};

  % Traiettoria
  \draw[traj] (-0.6,0.9) .. controls (1.2,1.4) and (3.4,2.4) .. (6.6,2.7);

  % Finestra sulla linea temporale
  \fill[gray!20, draw=black!50] (\tk,\yaxis+0.22) rectangle (\tkW,\yaxis-0.22);

  % Punti di ingresso/uscita su frontiera R
  \coordinate (EinB)  at (0,1.05);
  \coordinate (EoutB) at (6,2.62);
  \fill[red] (EinB) circle (1.6pt);
  \node[openmark] at (EoutB) {};
  \node[font=\scriptsize, anchor=south west, yshift=1em] at (EinB)  {$t^{\mathrm{in}}_{R}=1.4$};
  \node[font=\scriptsize, anchor=south east, yshift=-0.4em] at (EoutB) {$t^{\mathrm{out}}_{R}=3.2$};

  % Proiezioni su asse temporale
  \draw[proj] (EinB)  -- (\tinB,\yaxis);
  \draw[proj] (EoutB) -- (\toutB,\yaxis);

  % Asse temporale
  \draw[taxis] (-0.2,\yaxis) -- (6.6,\yaxis);
  \foreach \x in {0,...,6}{
    \draw (\x,\yaxis+0.08) -- (\x,\yaxis-0.08);
    \node[font=\scriptsize] at (\x,\yaxis-0.35) {\x.0};
  }

  % Marker ingresso/uscita su asse
  \node[tmark/in]  at (\tinB,\yaxis) {};
  \node[tmark/out] at (\toutB,\yaxis) {};
  
\end{tikzpicture}
\subcaption{Rotta B: \textbf{non copre} la finestra}
\end{subfigure}

\caption{Occupazione \(O_R(k)\) per la regione \(R\): conteggio dei voli la cui permanenza copre interamente la finestra \(I_k=[t_k,t_k{+}W)\).}
\label{fig:occupazione-subfig}
\end{figure}

L'implementazione algoritmica non valuta la condizione di copertura \eqref{eq:occ-finestra-copertura} per ogni volo e ogni finestra, ma sfrutta il principio dell'array di differenze. La proposizione seguente formalizza l’identità che adottiamo nell’implementazione.

\begin{proposition}[Equivalenza tra occupancy a copertura e a differenze]
\label{prop:occupancy-equivalence}
Sia \(R\) una regione di volo, \(I_k = [t_k, t_k+W)\) una finestra temporale su una griglia \(\mathcal{T}\) di passo \(\lambda>0\) con origine \(\tau_0\), e sia \(\mathcal{F}_R\) l'insieme dei voli che attraversano \(R\).
Definiamo innanzitutto l'\emph{occupancy a copertura} come
\begin{equation*}
  O_R^{\text{cop}}(k) \coloneqq 
  \#\Bigl\{\, f \in \mathcal{F}_R \;\Bigm|\; 
  t^{\mathrm{in}}_{f,R} \le t_k \ \wedge\ t^{\mathrm{out}}_{f,R} \ge t_k+W \Bigr\}.
\end{equation*}
Definiamo poi l'\emph{occupancy a differenze} come
\begin{equation*}
  O_R^{\text{diff}}(k) \coloneqq \sum_{i=0}^k D_R(i),
  \qquad \text{dove} \qquad
  D_R(i) \coloneqq \sum_{f \in \mathcal{F}_R} \delta_f(i),
\end{equation*}
e la funzione di contributo del singolo volo \(f\) come
\begin{equation*}
  \delta_f(i) \coloneqq
  \begin{cases}
    +1 & \text{se } i = k_s(f), \\
    -1 & \text{se } i = k_e(f) + 1, \\
    0  & \text{altrimenti}.
  \end{cases}
\end{equation*}

Gli indici \(k_s(f)\) e \(k_e(f)\) sono scelti in modo che il volo \(f\) contribuisca esattamente alle finestre che risultano completamente coperte dall'intervallo di permanenza in \(R\), ossia \([t^{\mathrm{in}}_{f,R}, t^{\mathrm{out}}_{f,R}]\).

Perché la finestra \(I_k = [t_k, t_k+W)\) sia interamente contenuta in \([t^{\mathrm{in}}_{f,R}, t^{\mathrm{out}}_{f,R}]\) devono valere le disuguaglianze
\begin{equation*}
  t^{\mathrm{in}}_{f,R} \le t_k 
  \qquad\text{e}\qquad 
  t_k + W \le t^{\mathrm{out}}_{f,R}.
\end{equation*}
Poiché i nodi della griglia sono definiti da
\begin{equation*}
  t_k = \tau_0 + k\,\lambda,
\end{equation*}
sostituendo otteniamo il sistema
\begin{equation*}
  t^{\mathrm{in}}_{f,R} \le \tau_0 + k\,\lambda,
  \qquad
  \tau_0 + k\,\lambda + W \le t^{\mathrm{out}}_{f,R}.
\end{equation*}
Risolvendo rispetto a \(k\):
\begin{align*}
  t^{\mathrm{in}}_{f,R} \le \tau_0 + k\,\lambda 
  &\iff k \ge \frac{t^{\mathrm{in}}_{f,R} - \tau_0}{\lambda}, \\[0.5em]
  \tau_0 + k\,\lambda + W \le t^{\mathrm{out}}_{f,R} 
  &\iff k \le \frac{t^{\mathrm{out}}_{f,R} - W - \tau_0}{\lambda}.
\end{align*}
Dunque gli indici interi \(k\) per cui il volo \(f\) copre completamente la finestra \(I_k\) sono esattamente quelli che soddisfano
\begin{equation*}
  \frac{t^{\mathrm{in}}_{f,R} - \tau_0}{\lambda} 
  \;\le\; k \;\le\;
  \frac{t^{\mathrm{out}}_{f,R} - W - \tau_0}{\lambda}.
\end{equation*}
Imponendo che \(k\) sia intero, otteniamo quindi
\begin{equation*}
  k_s(f) \coloneqq 
  \left\lceil \frac{t^{\mathrm{in}}_{f,R} - \tau_0}{\lambda} \right\rceil,
  \qquad
  k_e(f) \coloneqq 
  \left\lfloor \frac{t^{\mathrm{out}}_{f,R} - W - \tau_0}{\lambda} \right\rfloor.
\end{equation*}
Se \(k_s(f) > k_e(f)\) l'intervallo \([t^{\mathrm{in}}_{f,R}, t^{\mathrm{out}}_{f,R}]\) non contiene alcuna finestra \(I_k\) per intero e il volo \(f\) non contribuisce ad alcun indice; altrimenti contribuisce per tutti e soli gli indici interi \(k\) tali che \(k_s(f) \le k \le k_e(f)\). Con queste definizioni vale l'identità
\begin{equation}
\label{eq:occupancy-equivalence}
  O_R^{\text{cop}}(k) = O_R^{\text{diff}}(k) 
  \quad \forall k \ge 0.
\end{equation}
\end{proposition}

\begin{proof}
Fissiamo una regione $R$ e un indice $k$. Per definizione, $O_R^{\text{cop}}(k)$ è il numero di voli $f$ tali che l’intera finestra
\[
I_k = [t_k,\,t_k{+}W)
\]
è contenuta nell’intervallo di permanenza del volo in $R$, cioè in
\[
[t^{\mathrm{in}}_{f,R},\,t^{\mathrm{out}}_{f,R}] .
\]
Questo è equivalente a dire che esistono due indici interi $k_s(f)$ e $k_e(f)$ tali che il volo $f$ copre esattamente le finestre con
\[
k_s(f) \le k \le k_e(f),
\]
e non copre nessun’altra finestra.  
In altre parole,
\[
O_R^{\text{cop}}(k) 
= \#\bigl\{\, f \in \mathcal{F}_R \;\bigm|\; k_s(f) \le k \le k_e(f) \bigr\}.
\]
Per ogni volo $f$ poniamo
\[
  \delta_f(i) \coloneqq
  \begin{cases}
    +1 & \text{se } i = k_s(f), \\
    -1 & \text{se } i = k_e(f) + 1, \\
    0  & \text{altrimenti},
  \end{cases}
\]
e $D_R(i) = \displaystyle\sum_{f\in \mathcal{F}_R} \delta_f(i)$.  
L’occupancy a differenze è la somma cumulata:
\begin{equation*}
\label{eq:sum-cum}
  O_R^{\text{diff}}(k) \coloneqq \sum_{i=0}^k D_R(i).
\end{equation*}

Per capire che cosa stiamo contando, è sufficiente guardare il contributo di un singolo volo $f$.  
Definiamo
\[
S_f(k) \coloneqq \sum_{i=0}^k \delta_f(i).
\]
Poiché $\delta_f(i)$ è diverso da zero solo in $i=k_s(f)$ e in $i=k_e(f){+}1$, abbiamo tre casi:
\begin{itemize}
  \item se $k < k_s(f)$, non abbiamo ancora sommato il $+1$, quindi $S_f(k)=0$
  \item se $k_s(f) \le k \le k_e(f)$, abbiamo sommato il $+1$ in $k_s(f)$ ma non ancora il $-1$ in $k_e(f){+}1$, quindi $S_f(k)=1$
  \item se $k > k_e(f)$, nella somma sono entrati sia il $+1$ (in $k_s(f)$) sia il $-1$ (in $k_e(f){+}1$), che si annullano, quindi $S_f(k)=0$
\end{itemize}

Dunque $S_f(k)$ è semplicemente la funzione indicatrice del fatto che il volo $f$ copre la finestra $I_k$:
\[
S_f(k) =
\begin{cases}
1 & \text{se } k_s(f) \le k \le k_e(f),\\
0 & \text{altrimenti}.
\end{cases}
\]
\medskip
Osserviamo ora che, per definizione di $D_R(i)$,
\[
  D_R(i) = \sum_{f\in\mathcal{F}_R} \delta_f(i).
\]
Sommando fino all’indice $k$ otteniamo
\[
  O_R^{\text{diff}}(k)
  = \sum_{i=0}^k D_R(i)
  = \sum_{i=0}^k \sum_{f\in\mathcal{F}_R} \delta_f(i)
  = \sum_{f\in\mathcal{F}_R} \underbrace{\sum_{i=0}^k \delta_f(i)}_{S_f(k)}.
\]
Dunque
\[
  O_R^{\text{diff}}(k) = \sum_{f\in\mathcal{F}_R} S_f(k).
\]

Ma abbiamo appena visto che $S_f(k)$ vale $1$ se e solo se il volo $f$ copre la finestra $I_k$
(ossia se $k_s(f) \le k \le k_e(f)$), e vale $0$ altrimenti.
Quindi la somma
\[
  \sum_{f\in\mathcal{F}_R} S_f(k)
\]
non è altro che il numero di voli che coprono la finestra $I_k$, cioè
\[
  O_R^{\text{diff}}(k)
  = \#\bigl\{\, f \in \mathcal{F}_R \;\bigm|\; k_s(f) \le k \le k_e(f) \bigr\}.
\]
Confrontando con l’espressione trovata all’inizio per $O_R^{\text{cop}}(k)$ si ottiene, per ogni $k$,
\[
  O_R^{\text{cop}}(k) = O_R^{\text{diff}}(k),
\]
cioè l’identità~\eqref{eq:occupancy-equivalence}.
\end{proof}



 \paragraph{Esempio.}
Consideriamo l'intervallo \([0,100]\) e, per semplicità, poniamo \(W=\lambda=10\).
Le finestre semi–aperte sono \(I_k=[t_k,t_k{+}10)\) con \(t_k=k\cdot10\) e \(k=0,\dots,9\).
I voli (intervalli di permanenza in \(R\)) sono:
\[
F_1=[15,45],\quad
F_2=[25,35],\quad
F_3=[28,62],\quad
F_4=[5,10],\quad
F_5=[50,80].
\]

\medskip
\noindent\emph{Passo 1: indici discreti \(k_s\) e \(k_e\).}
Nel nostro caso \(\tau_0=0\) e \(\lambda=W=10\). Si ottiene:
\[
\begin{array}{c|cc|cc}
\text{Volo} & t^{\mathrm{in}} & t^{\mathrm{out}} & k_s=\lceil t^{\mathrm{in}}/10\rceil & k_e=\lfloor (t^{\mathrm{out}}-10)/10\rfloor\\\hline
F_1 & 15 & 45 & 2 & 3\\
F_2 & 25 & 35 & 3 & 2\\
F_3 & 28 & 62 & 3 & 5\\
F_4 & 5  & 10 & 1 & 0\\
F_5 & 50 & 80 & 5 & 7\\
\end{array}
\]
I voli \(F_2\) e \(F_4\) non contribuiscono poiché \(k_s>k_e\) (durata insufficiente a coprire interamente una finestra di ampiezza \(W\)).

\medskip
\noindent\emph{Passo 2: differenze su \(k\).}
Per ogni volo valido, aggiungiamo \(+1\) in \(k_s\) e \(-1\) in \(k_e{+}1\); quindi aggreghiamo per \(k\):
\[
\begin{array}{c|cccccccccc}
k & 0&1&2&3&4&5&6&7&8&9\\\hline
\delta_{F_1}(k) & 0&0&{+}1&0&{-}1&0&0&0&0&0\\
\delta_{F_3}(k) & 0&0&0&{+}1&0&0&{-}1&0&0&0\\
\delta_{F_5}(k) & 0&0&0&0&0&{+}1&0&0&{-}1&0\\\hline
D(k)=\sum_f \delta_f(k) & 0&0&{+}1&{+}1&{-}1&{+}1&{-}1&0&{-}1&0
\end{array}
\]

\medskip
\noindent\emph{Passo 3: somma cumulata.}
L'occupancy è la somma cumulata:
\[
\begin{array}{c|cccccccccc}
k & 0&1&2&3&4&5&6&7&8&9\\\hline
O(k) & 0&0&1&2&1&2&1&1&0&0
\end{array}
\]


\clearpage
\section{Progettazione degli algoritmi di calcolo delle metriche}
In questa sezione si descrive la progettazione degli algoritmi che permettono di calcolare le serie temporali di \emph{throughput} e \emph{occupancy} per ciascuna regione.  
Gli algoritmi sono presentati come pseudocodice e costituiscono il riferimento concettuale per l'implementazione software sviluppata successivamente.

\subsection{Throughput}
\label{sec:progettazione-metriche}
Dal punto di vista algoritmico, il calcolo del \emph{throughput} introdotto in Sezione~\ref{subsec:throughput-teoria} parte da una tabella di eventi \(E\), in cui ogni riga rappresenta l’ingresso di un volo in una regione \(R\) al tempo \(t^{\mathrm{in}}_{f,R}\).  
L’obiettivo è ottenere, per ogni regione e per ogni finestra temporale \(I_k = [t_k, t_k+W)\) della griglia \(\mathcal{T}\), il numero di ingressi che cadono all’interno di \(I_k\).

Per evitare di scorrere tutti gli eventi per ogni finestra (approccio inefficiente su grandi volumi di dati), si sfrutta l’identità della Proposizione~\ref{prop:two-cumulatives}, che esprime il throughput come differenza di due valori di una funzione cumulata \(N_R(t)\), dove \(N_R(t)\) è il numero di ingressi avvenuti in \(R\) prima dell’istante \(t\). In particolare,
\[
  T_R(k) \;=\; N_R(t_k{+}W) - N_R(t_k),
\]
così che per ogni finestra siano sufficienti due sole valutazioni della cumulata.
L’Algoritmo~\ref{alg:throughput-anyW-naturale} traduce questa idea come segue.  
Anzitutto costruisce la griglia temporale tra \(\tau_0\) e \(\tau_1\) con passo \(\lambda\), determinando il numero di nodi \(K\) e quindi le finestre \(I_k\) (riga~1).  
Successivamente filtra gli eventi, conservando soltanto gli ingressi con tempo in \([\tau_0,\,\tau_1{+}W)\), così da coprire anche le finestre che iniziano vicino al bordo superiore dell’orizzonte (riga~2), e per ogni regione \(R\) aggrega gli eventi che condividono lo stesso timestamp \(t\), ottenendo i conteggi elementari \(n_R(t)\) (riga~3).  
Su questa base, per ciascuna regione costruisce la funzione cumulata “stretta” \(N_R(t)\), che per ogni istante \(t\) restituisce quanti ingressi sono già avvenuti prima di \(t\) (riga~4).  
Infine, per ogni regione e per ogni finestra \(I_k\), valuta i valori della cumulata al bordo sinistro e al bordo destro e calcola il throughput come differenza (righe~7–9), producendo la tabella di output \(T\) ordinata per \((R,t_k)\).

Questo schema è indipendente dalla tecnologia utilizzata e si presta a un’implementazione vettoriale sulle strutture tabellari del dataset, riducendo il costo computazionale rispetto a una scansione finestra per finestra degli eventi grezzi.
\clearpage

\begin{algorithm}[ht]
  \caption{Throughput per regione.}
  \label{alg:throughput-anyW-naturale}
  \begin{algorithmic}[1]
    \Require tabella eventi $E$ con colonne $R$ (regione) e $t^{\mathrm{in}}$; estremi $\tau_0<\tau_1$; passo $\lambda$; ampiezza $W>0$.
    \Ensure tabella $T$ con righe $(R,\,t_k,\,t_k+W,\,T_R(k))$ ordinate per $(R,\,t_k)$.
    \EnsureEndSkip{}
    \State \textbf{Prepara la griglia.} Calcola quanti passi stanno tra $\tau_0$ e $\tau_1$.
    \State \textbf{Tieni solo gli eventi che servono.} Filtra $E$ conservando gli ingressi con tempo nell’intervallo $[\tau_0,\,\tau_1{+}W)$: così copri anche le finestre che iniziano vicino a $\tau_1$.
    \State \textbf{Conta per istante.} Raggruppa gli eventi che hanno la stessa regione $R$ e lo stesso timestamp $t$; assegna a ciascuna coppia $(R,t)$ il conteggio $n_R(t)$, pari al numero di ingressi occorsi \emph{esattamente} a $t$ in $R$.
    \State \textbf{Costruisci la cumulata “stretta”.} Per ciascuna regione ordina i tempi e costruisci $N_R(t)$ come “quanti ingressi sono già avvenuti prima di $t$”.
    \State \textbf{Inizializza l’output.} Metti $T \gets [\ ]$.
    \For{ogni regione $R$}
      \For{$k \gets 0$ \textbf{to} $K{-}1$}
        \State \textbf{Valuta la finestra.} Prendi il valore di cumulata al bordo sinistro ($N_R(t_k)$) e al bordo destro ($N_R(t_k{+}W)$). Il primo dice “quanti ingressi erano già avvenuti prima che iniziasse la finestra”, il secondo “quanti ce ne sono prima che la finestra finisca”.
        \State \textbf{Calcola il throughput:}
            \[
                T_R(k) \;\gets\; N_R(t_k{+}W) - N_R(t_k)
            \]
        \State \textbf{Accoda il risultato.} Aggiungi a $T$ la riga $(R,\,t_k,\,t_k{+}W,\,T_R(k))$.
      \EndFor
    \EndFor
    \State Ordina $T$ per $(R,\,t_k)$
    \Return $T$.
  \end{algorithmic}
\end{algorithm}
\clearpage
\subsection{Occupancy}
\label{subsec:occupancy-impl}
Dal punto di vista algoritmico, il calcolo dell’\emph{occupancy} introdotto in Sezione~\ref{subsec:occupancy-teoria} parte da una tabella di attraversamenti \(D\), in cui ogni riga rappresenta un volo \(f\) che entra in una regione \(R\) al tempo \(t^{\mathrm{in}}_{f,R}\) ed esce al tempo \(t^{\mathrm{out}}_{f,R}\).  
L’obiettivo è determinare, per ogni regione e per ogni finestra temporale \(I_k = [t_k, t_k+W)\) della griglia \(\mathcal{T}\), il valore \(O_R(k)\) che conta i voli presenti in \(R\) per l’intera durata della finestra. 
L’algoritmo lavora sugli indici discreti della griglia. Per ogni volo filtrato determina:
\begin{itemize}
  \item l’indice \(k_s\) della prima finestra che il volo contribuisce a occupare
  \item l’indice \(k_e\) dell’ultima finestra da considerare
\end{itemize}
Invece di aggiornare direttamente \(O_R(k)\) per tutti gli indici \(k_s \le k \le k_e\), registra soltanto due variazioni: un incremento che fa crescere l’occupancy a partire da \(k_s\), e un decremento che la fa diminuire dopo l’ultima finestra influenzata dal volo.  
Aggregando queste variazioni per coppia \((R,k)\) si ottiene, per ciascuna regione, un array di differenze \(D_R(k)\); calcolandone la somma cumulata rispetto a \(k\) si ricostruiscono i valori \(O_R(k)\).

L’Algoritmo~\ref{alg:occupancy-anyW-diff} implementa questo schema: prepara la griglia temporale e filtra i voli che non possono contribuire ad alcuna finestra dell’orizzonte considerato (righe~1–2), costruisce l’array di differenze aggiungendo per ogni volo gli eventi di ingresso e uscita (righe~4–10), aggrega le variazioni per regione e indice \(k\) (righe~11–12) e, infine, per ciascuna regione, scorre gli indici in ordine crescente accumulando le differenze per ottenere l’occupancy \(O_R(k)\) e popolare la tabella di output \(O\) (righe~14–20).
\clearpage

\begin{algorithm}[H]
  \caption{Occupancy per regione.}
  \label{alg:occupancy-anyW-diff}
    \vspace{-1.5em} %fix grafico dovuto a un buh
  \begin{algorithmic}[1]
    \Require tabella eventi $D$ con colonne $R$ (regione), $t^{\mathrm{in}}$, $t^{\mathrm{out}}$; estremi $\tau_0<\tau_1$; passo $\lambda$; ampiezza $W>0$.
    \Ensure Tabella $O$ con righe $(R,\,t_k,\,t_k+W,\,O_R(k))$ ordinate per $(R,\,t_k)$.
    \EnsureEndSkip{}
    \State \textbf{Preparazione} Calcola gli istanti $t_k = \tau_0 + k \lambda$ per $k \in [0, K{-}1]$.
    \State \textbf{Filtraggio dei voli.} Conserva solo i voli $f$ con $t^{\mathrm{out}}_{f,R} \ge \tau_0+W$ e $t^{\mathrm{in}}_{f,R} \le t_{K-1}$.
    \State Inizializza l'array di differenze $E \gets [\ ]$.
    \Comment{E: $(\mathbf{R}, \mathbf{k}, \mathbf{\delta})$}
    \For{ogni volo $f$ filtrato}
        \State \quad $k_{s} \gets \max\left(0, \left\lceil \frac{t^{\mathrm{in}}_{f,R} - \tau_0}{\lambda} \right\rceil\right)$
        \State \quad $k_{e} \gets \min\left(K{-}1, \left\lfloor \frac{t^{\mathrm{out}}_{f,R} - W - \tau_0}{\lambda} \right\rfloor\right)$
        \If{$k_{s} \le k_{e}$}
            \State Registra evento \textbf{ingresso}: $E \gets E \cup \{(R, k_{s}, +1)\}$.
            \State Registra evento \textbf{uscita}: $E \gets E \cup \{(R, k_{e} + 1, -1)\}$.
        \EndIf
    \EndFor
    \State $B \gets \text{Aggrega } E \text{ sommando } \mathbf{\delta} \text{ per } (R, \mathbf{k}) \text{ (ottenendo } \mathbf{\delta_{sum}}).$
    \State Estendi $\mathbf{B}$ per includere tutte le coppie $(R, \mathbf{k})$ della griglia, ponendo $\mathbf{\delta} \gets 0$ per gli istanti senza evento.
  
    \State Ordina $B$ per $(R,k)$.
    \For{ogni $R$}
      \State $s \gets 0$
      \For{$k = 0,\dots,K{-}1$} \Comment scansiona gli indici in ordine crescente
        \State $s \gets s + \bigl(\delta(R,k))\ $
        \State $O_R(k) \gets s$
      \EndFor
      \State \textbf{Accoda il risultato.} Aggiungi a $O$ la riga $(R,\,t_k,\,t_k{+}W,\,O_R(k))$.
    \EndFor
\State \Return $O$
   
  \end{algorithmic}
\end{algorithm}